\chapter{The model}
\label{chap:model}

Space is $\R^3$ with $\up$ pointing up; $\ang(a,b)\in[0,\pi]$ is the unoriented
angle between vectors. Time is discrete with period $\delta>0$, which is both the
control period and the camera frame period. Visibility is required at frame
instants. Every statement in this blueprint is about the full rigid-body
quadrotor below.

\begin{definition}[Stepping]
\label{def:steps}
\lean{VisionTracking.StepsWith}
\leanok
Euler update: $x_{k+1}=x_k+\delta v_k$. A position-only observer therefore recovers
velocities from consecutive positions.
\end{definition}

\section{Target}

\begin{definition}[Target parameters]
\label{def:target-params}
\lean{VisionTracking.TargetParams}
\leanok
Speed band $0<v_{\min}\le v_{\max}$, minimum turn radius $R>0$, maximum along-track
acceleration $a_\parallel\ge0$, maximum jerk $j\ge0$, maximum climb/dive angle
$\gamma<\pi/2$, and detection size $s>0$.
\end{definition}

\begin{definition}[Admissible target trajectory]
\label{def:target-traj}
\lean{VisionTracking.TargetTraj}
\leanok
\uses{def:steps, def:target-params}
The speed stays in the band. The velocity increment splits into an along-track
part (at most $\delta a_\parallel$) and a normal part (at most
$\delta\norm{v}^2/R$, a turn-radius limit in every direction). The acceleration
changes by at most $\delta j$ per step, and the climb angle stays within $\gamma$.
\end{definition}

\section{Camera}

\begin{definition}[Camera parameters]
\label{def:camera-params}
\lean{VisionTracking.CameraParams}
\leanok
Half-angle $\alpha$ of the inscribed field-of-view cone; horizontal field of view
$\Phi$ over $W$ pixels; detection threshold $n$ pixels; hinge range
$[\eta_{\mathrm{lo}},\eta_{\mathrm{hi}}]\subseteq[-\pi/2,\pi/2]$ about the body
axis $b_2$; hinge rate $\omega_\eta$. A fixed camera has
$\eta_{\mathrm{lo}}=\eta_{\mathrm{hi}}$.
\end{definition}

\begin{definition}[Detection range]
\label{def:detect-range}
\lean{VisionTracking.CameraParams.detectRange}
\leanok
\uses{def:camera-params}
$\ddet(s)=s/(2\tan(n\Phi/2W))\approx sW/(n\Phi)$.
\end{definition}

\begin{definition}[Yaw-parametrised boresight]
\label{def:cam-axis}
\lean{VisionTracking.camAxis, VisionTracking.bodyForward, VisionTracking.heading}
\leanok
For a thrust axis $b_3$ and a yaw $\psi$, $b_1$ is the heading projected onto
$b_3^\perp$ and normalised, and the boresight is $\cos\eta\,b_1+\sin\eta\,b_3$.
This parametrisation is used for the pointing geometry. Lemma~\ref{lem:yaw-covers}
shows that it reaches every body frame with a given thrust axis.
\end{definition}

\section{Drone: a rigid-body quadrotor}

\begin{definition}[Quadrotor parameters]
\label{def:quad-params}
\lean{VisionTracking.QuadParams}
\leanok
Gravity $g$; mass $m$; principal inertias $J_x,J_y,J_z$; arm length $L$; maximum
thrust per motor $f_{\mathrm{motor}}$ with $mg<4f_{\mathrm{motor}}$; yaw reaction
coefficient $k$; quadratic drag coefficient $c_D>0$; tilt envelope
$\theta_{\max}<\pi/2$; body-rate limit $\omega_{\max}$; and the speed envelope on
which the drag model is used.
\end{definition}

\begin{definition}[Motor mixing, ``+'' layout]
\label{def:mixing}
\lean{VisionTracking.totalThrust, VisionTracking.QuadParams.torque}
\leanok
Motors $0,1,2,3$ sit on $+b_1,+b_2,-b_1,-b_2$. The collective thrust is
$F=\sum_i f_i$ and the body torque is
$\tau=\bigl(L(f_1-f_3),\ L(f_2-f_0),\ k(f_0-f_1+f_2-f_3)\bigr)$.
\end{definition}

\begin{definition}[Cross product, Rodrigues rotation, inertia]
\label{def:rotate}
\lean{VisionTracking.cross, VisionTracking.rotate, VisionTracking.worldRate,
  VisionTracking.QuadParams.inertia, VisionTracking.QuadParams.inertiaInv,
  VisionTracking.frameCam}
\leanok
$\mathrm{rotate}(\omega)$ turns vectors by the angle $\theta=\norm{\omega}$ about
$\hat\omega$:
$v\mapsto\cos\theta\,v+\sin\theta\,(\hat\omega\times v)+(1-\cos\theta)\inner{\hat\omega}{v}\hat\omega$.
The camera boresight in the body frame is $\cos\eta\,b_1+\sin\eta\,b_3$.
\end{definition}

\begin{definition}[Admissible quadrotor trajectory]
\label{def:quad-traj}
\lean{VisionTracking.QuadTraj, VisionTracking.QuadTraj.cam}
\leanok
\uses{def:steps, def:quad-params, def:mixing, def:rotate, def:camera-params}
The state is position $x$, velocity $v$, a right-handed orthonormal body frame
$(b_1,b_2,b_3=b_1\times b_2)$ (which encodes roll, pitch and yaw), body rates
$\Omega$, motor thrusts $f_0,\dots,f_3$, and hinge angle $\eta$. It evolves by
\[
  v_{k+1}=v_k+\delta\Bigl(\tfrac{F_k}{m}b_{3,k}-g\up-c_D\norm{v_k}v_k\Bigr),\qquad
  \Omega_{k+1}=\Omega_k+\delta J^{-1}\bigl(\tau_k-\Omega_k\times J\Omega_k\bigr),
\]
and each $b_i$ is advanced by $\mathrm{rotate}(\delta\,\omega_k)$, where $\omega_k$
is the world-frame angular velocity. Constraints: $0\le f_i\le f_{\mathrm{motor}}$,
tilt at most $\theta_{\max}$, $\norm{\Omega}\le\omega_{\max}$, speed within the
envelope, and the hinge range and rate.
\end{definition}

\section{Visibility and strategies}

\begin{definition}[Visibility]
\label{def:visible}
\lean{VisionTracking.Visible, VisionTracking.KeepsInView}
\leanok
\uses{def:detect-range, def:quad-traj, def:target-traj}
The target at $p$ is visible from a camera at $x$ with boresight $c$ if
$\norm{p-x}\le\ddet(s)$ and $\ang(c,p-x)\le\alpha$.
\end{definition}

\begin{definition}[Strategies]
\label{def:strategy}
\lean{VisionTracking.Strategy, VisionTracking.PosStrategy, VisionTracking.QuadTraj.state}
\leanok
\uses{def:quad-traj, def:target-traj}
A strategy maps target trajectories to quadrotor trajectories causally: the motor
commands at step $k$, and so the state at step $k+1$, depend on the target only
through steps $0,\dots,k$. A position-only strategy sees only target positions,
which bearing plus apparent size determine for a target of known size.
\end{definition}

\begin{definition}[Maintainable]
\label{def:maintainable}
\lean{VisionTracking.Maintainable, VisionTracking.MaintainableFullInfo,
  VisionTracking.TargetTraj.StartsAt, VisionTracking.Maintainable.fullInfo}
\leanok
\uses{def:strategy, def:visible}
Tracking is maintainable from a target start $(p_0,v_0,a_0)$ if some position-only
strategy keeps every admissible target with that start in view at every frame.
Impossibility results are proved against full-information strategies, which makes
them stronger.
\end{definition}
